\scanpage{321}
and we have
\[
(38)\qquad (-\delta\unsure{v|\nu})B=ac+bc-\unsure{v|\nu}bd
\]
which gives $\uline{u}$ again in terms of $a,b,c,d$, \uline{thanks to} the formula $\uline{\overline{\mathrm{IV}}}(48)$ (page $453'$)
\[
(39)\qquad \uline{u}=\begin{pmatrix}\dfrac1D&\inserted{B/D}\deleted{\unsure{\underline u=}}\\0&1\end{pmatrix}
=\begin{pmatrix}
\dfrac{-\unsure{v|\nu}(ad-bc)}{-\unsure{v|\nu}d^2+cd+d^2}&\dfrac{-ac+bc-\unsure{v|\nu}bd}{-\unsure{v|\nu}\unsure{d|a}^2+cd+d^2}\\0&1
\end{pmatrix}.
\]

\uline{Remark}. On the other hand, the datum of $\beta$ \inserted{\uline{alone}} \inserted{(without the ``sign'' $\varepsilon$)} does not allow one to recover $\mu$. But the datum of $(\beta,\unsure{\varepsilon})$ does, for from it one deduces the knowledge of
\[
(\unsure{40})\qquad \eta=\varepsilon\beta\sigma(\beta)^{-1}\overset{\mathrm{dfn}}=\begin{pmatrix}R&S\\T&U\end{pmatrix},
\]
which will be related to $\zeta=\alpha\rho(\alpha)^{-1}$ by
\[
\zeta=\eta\varepsilon_1(\mu-1)
\]
and by the theory \unsure{of the|of} (II) we will have
\[
(41)\qquad \mu=1/U
\]
(but in terms of $\beta\sigma(\beta)^{-1}\overset{\mathrm{dfn}}=\begin{pmatrix}R_0&S_0\\T_0&U_0\end{pmatrix}$,
we will have \deleted{\illegible{1}} \unsure{)}
\[
(42)\qquad \mu=\frac{\varepsilon}{U_0}\qquad)
\]
i.e. it is \inserted{a priori} not known from $U_0$ alone, but only via $U_0$ and $\varepsilon$. \deleted{\unsure{R}}

To make this point precise, we are going to look at \inserted{when and} how, for fixed $(\beta,\varepsilon)$, one can complete with $\alpha,\mu$ so as to have a
\mspage{320}

\scanpage{322}
\mspage{461}
section $(\alpha,\beta,\mu,\varepsilon)$ of $\mathcal G^{*}_{\rho,\sigma}$. A necessary condition is that $U_0$ be invertible (it depends only on $\beta$, not on $\varepsilon$) --- and then one can choose $\varepsilon$ arbitrarily in $\Gamma\mu_2$, \deleted{and}\inserted{$\varepsilon\in\Gamma\mu_2$} $\eta=\varepsilon\beta\sigma(\beta)^{-1}$ will then be a section of the \inserted{\unsure{subscheme}} $A_\sigma$ \deleted{\illegible{1}} of $G$ of solutions of $\det\eta=1$, $\operatorname{Tr}\eta\sigma=\operatorname{Tr}\sigma$ (cf. lemma page 455), hence by the theory of (II), for \deleted{\unsure{completing taking}} $\mu=1/U=\varepsilon/U_0$, \deleted{\illegible{2}} and defining
\[
\zeta=\eta\varepsilon_1(\mu-1),
\]
$\zeta$ will indeed be expressible \inserted{locally} in the form $\alpha\rho(\alpha)^{-1}$, the indeterminacy being moreover in $T_\rho$ via $\alpha\longmapsto\alpha'=\alpha\uline b$, $\uline b\in\Gamma T_\rho$.

This does show that $\mu$ can \unsure{thereof|not} take, for a fixed $\beta$, all the values $\varepsilon/U_0$, where $\beta\sigma(\beta)^{-1}=\begin{pmatrix}R_0&\unsure{S_0|C_0}\\T_0&U_0\end{pmatrix}$ \ldots
\mspage{321}

\scanpage{323}
\mspage{462}
($\uline{\overline{\mathrm{VI}}}$) Characterization \unsure{of the|of the} modular triple $(\varepsilon_1,\rho,\sigma)$.

$\unsure{1^\circ})$ Let $G,SG,PG$ be associated forms of $\mathrm{Gl}(2),\mathrm{Sl}(2),\mathrm{GP}(1)$ over a base $S$ --- thus corresponding: a torsor under $\mathrm{GP}(1)_S$, or again: a bundle of projective lines \inserted{$P$} over $S$, or an Azumaya Algebra $M$ of rank $\unsure{1|4}$, and let us consider two sections \inserted{\illegible{1}} $\widetilde\rho,\widetilde\sigma$ of $PG$, such that
\[
(1)\qquad \unsure{a)}\quad\widetilde\varepsilon_1=\widetilde\rho\widetilde\sigma^{-1}
\]
be \deleted{\illegible{1}} a \uline{regular} unipotent section, and that the section
\[
(2)\qquad \widetilde\varepsilon_0\overset{\mathrm{dfn}}=\widetilde\rho^{-1}(\widetilde\varepsilon_1)=\widetilde\sigma^{-1}(\widetilde\varepsilon_1)
\]
which is deduced from it be \deleted{\unsure{everywhere}} \inserted{in each fiber} inequivalent to $\widetilde\varepsilon_0$ (i.e. that the sections $x_0,x_1$ of $P$ which correspond to $\varepsilon_0,\varepsilon_1$ be \deleted{\unsure{everywhere}} distinct in each fiber).

\margin{\uline{N.B.} If $\varepsilon_0,\varepsilon_1$ are the canonical liftings of $\widetilde\varepsilon_0,\widetilde\varepsilon_1$ into $\unsure{G|P}$, \unsure{regular}, \inserted{with $\unsure{\widetilde\sigma(\varepsilon_0)=\widetilde\rho(\varepsilon_0)=\varepsilon_1|\widetilde\sigma(\widetilde\varepsilon_0)=\widetilde\rho(\widetilde\varepsilon_0)=\widetilde\varepsilon_1}$} one can lift $\widetilde\rho,\widetilde\sigma$ in a \uline{unique} way to $\rho,\sigma$ so as to satisfy
$-\sigma^{-1}\rho=\varepsilon_0\quad(\Longleftrightarrow-\rho\sigma^{-1}=\varepsilon_1)$,
$(\rho+\sigma)^2=(\rho+\sigma)$.}

Then there exists a \uline{unique}$^*$ pinning
\[
(3)\qquad\deleted{\unsure{G_0=}}\quad\mathrm{Gl}(2)_S\xrightarrow[\sim]{\varphi}G
\]
and sections $\unsure{v|\nu},\unsure{w|\omega},\unsure{v|\nu}',\unsure{w|\omega}'$ of $\mathcal O_S$, satisfying
\[
(4)\qquad \unsure{v|\nu}+\unsure{v|\nu}'=0,\qquad \unsure{w|\omega}+\unsure{w|\omega}'=1,\qquad \unsure{v|\nu}\in\Gamma\mathcal O_S^*
\]
such that $\widetilde\rho,\widetilde\sigma$ are the images, under the composite $\mathrm{Gl}(2)_S\xrightarrow[\sim]{\varphi}G\longrightarrow PG$, of the sections
\[
(5)\qquad \rho=\begin{pmatrix}0&1\\\unsure{v|\nu}&\unsure{w|\omega}\end{pmatrix},\qquad
\sigma=\begin{pmatrix}0&-1\\\unsure{v|\nu}'&\unsure{w|\omega}'\end{pmatrix}
\]
\margin{\uline{N.B.}$^*$ Uniqueness \inserted{of $\unsure{v|\nu},\unsure{w|\omega}$ (i.e. of $\rho_0,\sigma_0$)} \unsure{is} not true without some restrictions on $S$, it suffices that every section $\tau$ of $\mathcal O_S$ such that $2\tau=0,\ \tau^2=0$ be zero.}
\mspage{322}

\scanpage{324}
of $\mathrm{Gl}(2)_S$ --- and conversely, for every pinning $\varphi$ and \unsure{every system} \deleted{\illegible{4}}\inserted{$\unsure{v|\nu},\unsure{w|\omega},\unsure{v|\nu}',\unsure{w|\omega}'$} satisfying (4) --- i.e.
\[
(4\mathrm{bis})\qquad \unsure{v|\nu}'=-\unsure{v|\nu},\quad \unsure{w|\omega}'=1-\unsure{w|\omega},\quad\begin{gathered}\unsure{v|\nu}\in\Gamma(S,\mathcal O_S^*)\\\unsure{w|\omega}\in\Gamma(S,\mathcal O_S)\end{gathered},
\]
the images \inserted{$\widetilde\rho,\widetilde\sigma$} under the composite $\mathrm{Gl}(2)_S\longrightarrow PG\simeq G/\mathbf G_m$ \deleted{satisfy}\inserted{of the} sections $\rho,\sigma$ of $\mathrm{Gl}(2)_S$ given by (5) satisfy the conditions above ($\widetilde\rho\widetilde\sigma^{-1}\overset{\mathrm{dfn}}=\widetilde\varepsilon_1$ is \uline{regular} unipotent, and everywhere inequivalent to $\widetilde\sigma^{-1}(\widetilde\varepsilon_1)=\widetilde\rho^{-1}(\widetilde\varepsilon_1)\overset{\mathrm{dfn}}=\widetilde\varepsilon_0$).

Setting then
\[
(6)\qquad \varepsilon_{\unsure{1|0}}=\begin{pmatrix}1&0\\1&1\end{pmatrix},\qquad \varepsilon_0=\begin{pmatrix}1&1/\unsure{v|\nu}\\0&1\end{pmatrix}
\]
we will have
\[
(7)\qquad \rho\sigma^{-1}=-\varepsilon_1,\qquad\sigma^{-1}\rho=-\varepsilon_0
\]
\[
(7\mathrm{bis})\qquad \varepsilon_0=\rho^{-1}(\varepsilon_1)=\sigma^{-1}(\varepsilon_1),\qquad\varepsilon_1=\rho(\varepsilon_0)=\sigma(\varepsilon_0)
\]
\uline{whence} the images under $\mathrm{Gl}(2)_S\longrightarrow PG$ of \deleted{$\rho$ and $\sigma$} $\varepsilon_1,\varepsilon_0$ are $\widetilde\varepsilon_1,\widetilde\varepsilon_0$. We have the formulas
\[
(8)\qquad \operatorname{Tr}\sigma=\unsure{w|\omega}',\quad\operatorname{Tr}\rho=\unsure{w|\omega},\quad\operatorname{Tr}(\sigma+\rho)=1
\]
\[
(9)\qquad \sigma+\rho=\begin{pmatrix}0&0\\0&1\end{pmatrix}\overset{\mathrm{dfn}}=N\deleted{\unsure{0}}
\]
\[
(10)\qquad \operatorname{Tr}\zeta N=D\quad\text{if}\quad\zeta=\begin{pmatrix}A&B\\C&\uline D\end{pmatrix},
\]
\[
(11)\qquad \det\sigma=\det\rho=-\unsure{v|\nu}=\unsure{-v'|v'}.
\]
\mspage{323}

\scanpage{325}
\mspage{463}
$2^\circ)$ The following conditions are equivalent
\diagram{\[
(12)\quad\left\{
\begin{tikzcd}[ampersand replacement=\&,column sep=large,row sep=large]
\boxed{\operatorname{Tr}\sigma=0}\arrow[r,Leftrightarrow]\arrow[d,Leftrightarrow]\&\operatorname{Tr}\rho=1\arrow[r,Leftrightarrow]\arrow[d,Leftrightarrow]\&\boxed{\widetilde\sigma^2=1}\arrow[d,Leftrightarrow]\\
\deleted{\mathrm{Tr}}\boxed{\unsure{w|\omega}'=0}\&\boxed{\unsure{w|\omega}=1}\&\sigma^2\in\Gamma\mathbf G_m\arrow[d,Leftrightarrow]\\
\&\&\boxed{\sigma^2=\unsure{v|\nu}\,\mathrm{id}.}
\end{tikzcd}\right.
\]}
They are also equivalent to this: denoting by $N'_\sigma$ the subgroup scheme of $\mathrm{Gl}(2)$ defined by the condition
\[
(13)\qquad \tau(\widetilde\sigma)=\widetilde\sigma\quad\text{i.e.}\quad\tau(\sigma)=\varepsilon\sigma,\quad\varepsilon\in\Gamma\mathbf G_m
\]
(we will necessarily have $\varepsilon^2=1$ i.e. $\varepsilon\in\Gamma\mu_2$)
the morphism
\[
(14)\qquad\begin{aligned}N'_\sigma&\longrightarrow\mu_2\\\tau&\longmapsto\varepsilon\end{aligned}
\]
is an \uline{epimorphism} --- so that we will have an exact sequence
\[
(15)\qquad1\longrightarrow T_\sigma\longrightarrow N'_\sigma\longrightarrow\mu_2\longrightarrow1,
\]
where \deleted{\unsure{Let us set}}
\[
(16)\qquad T_\sigma=\operatorname{Centr}_{\mathrm{Gl}(2)_S}(\sigma)=\{t\sigma+u\mid\det(t\sigma+u)\ \text{invertible}\}
\]
\[
(17)\qquad\text{We have}\quad N'_\sigma\subset N_\sigma\overset{\mathrm{dfn}}=\operatorname{Norm}_{\mathrm{Gl}(2)}(T_\sigma).
\]
If $2$ is a regular section of $\mathcal O_S$, then every section of $N_\sigma$ \inserted{over $S$} is a section of $N'_\sigma$ --- thus \unsure{if|in} \inserted{the} residue characteristics \unsure{are|is} $\ne2$ i.e. $2$ invertible, we will have
\[
(18)\qquad N'_\sigma=N_\sigma\qquad(\text{residue characteristics }\ne2).
\]
The condition under consideration means that $\sigma,\rho$ take the form
\mspage{324}

\scanpage{326}
\[
(19)\qquad\sigma=\begin{pmatrix}0&-1\\\underset{\substack{\Vert\\-\unsure{v|\nu}}}{\unsure{v|\nu}'}&0\end{pmatrix},\qquad\rho=\begin{pmatrix}0&1\\\unsure{v|\nu}&1\end{pmatrix}
\]
with $\unsure{v|\nu}+\unsure{v|\nu}'=0$.

These results of $1^\circ),2^\circ)$ are only a recapitulation of part of the results of the preceding sections $\uline{\overline{\mathrm{IV}}},\uline{\overline{\mathrm{V}}}$. \deleted{\unsure{Here}}, we have just studied conditions which amount to the determination \inserted{\unsure{the value of|the values of}} $\unsure{w|\omega}$ (or again \inserted{of} $\unsure{w|\omega}'$) in a \deleted{quite} explicit way. We are going to \unsure{turn to} conditions which \unsure{are to|will} serve to fix $\unsure{v|\nu}$ (or again $\unsure{v|\nu}'$) in an explicit way.

$3^\circ)$ Equivalent conditions
\[
(20)\quad\left\{\begin{aligned}
\det\sigma=1&\Longleftrightarrow\det\rho=1\Longleftrightarrow \unsure{v|\nu}=-1\Longleftrightarrow \unsure{v|\nu}'=1\\
&\Longleftrightarrow\varepsilon_0=\begin{pmatrix}1&-1\\0&1\end{pmatrix}.
\end{aligned}\right.
\]
These equivalences follow at once from (6) and (11)\unsure{.|!} The conditions also mean that $\rho$ and $\sigma$ have the form
\[
(21)\qquad \deleted{\rho}\inserted{\sigma}=\begin{pmatrix}0&-1\\1&\unsure{w|\omega}'\end{pmatrix},\qquad\rho=\begin{pmatrix}0&1\\-1&\unsure{w|\omega}\end{pmatrix}
\]
\[
(\unsure{w|\omega}+\unsure{w|\omega}'=1)
\]
$4^\circ)$ Conjunction of $2^\circ)$ and $3^\circ)$.

Equivalent conditions:
\mspage{325}

\scanpage{327}
\mspage{464}

\[
(22)\quad
\left\{
\begin{array}{ccccc}
\begin{gathered}
\text{Conjunction of the}\\
\text{conditions of }2^\circ\\
\text{and of the conditions}\\
\text{of }3^\circ
\end{gathered}
&\Longleftrightarrow&
\begin{gathered}\boxed{\sigma^2=-1}\\
\text{i.e. }\sigma^2+1=0\end{gathered}
&\Longleftrightarrow&
\boxed{\widetilde\sigma^2=\widetilde\rho^3=1}
\\
\Updownarrow&&&&\\
\unsure{v|\nu}=-1,\unsure{w|\omega}=1&\Longleftrightarrow&
\rho=\begin{pmatrix}0&1\\-1&1\end{pmatrix}
&\Longleftrightarrow&\operatorname{Tr}\rho=\det\rho=1
\\
\Updownarrow&&&&\\
\unsure{v|\nu}'=1,\unsure{w|\omega}'=0&\Longleftrightarrow&
\sigma=\begin{pmatrix}0&-1\\1&0\end{pmatrix}
&\Longleftrightarrow&(\operatorname{Tr}\sigma=0,\det\sigma=1)
\end{array}\right.
\]
\margin{\diagram{\begin{tikzcd}[
ampersand replacement=\&,
column sep=large,
row sep=small,
execute at end picture={\draw[-] (quadratic.north east) -- (quadratic.north west) -- (quadratic.south west) -- (quadratic.south east);}
]
\boxed{\widetilde\sigma^2=\widetilde\rho^3=1}
\arrow[r,Leftrightarrow]
\arrow[dr,Leftrightarrow]
\& \deleted{\boxed{\rho^3=-1}}
\\
\& |[alias=quadratic,inner sep=3pt]| {\rho^2-\rho+1=\unsure{\illegible{1}}}
\end{tikzcd}}
cf. corollary~2 below}

\uline{Proof.} The fact that the conjunction of the conditions of $2^\circ)$ and $3^\circ)$ is equivalent to \unsure{each} of the conditions of the $2^{\mathrm e}$ and $3^{\mathrm e}$ rows, is trivial. The equivalence with $\sigma^2=-1$ comes from the fact that \deleted{\illegible{1}} the condition $2^\circ$ means that \inserted{$\sigma^2$ \unsure{be|is} scalar, i.e.} $\sigma^2=\unsure{v|\nu}\,\mathrm{id}$, \deleted{\illegible{2}} whereas $3^\circ$ means $\unsure{v|\nu}=-1$, so $\sigma^2=-1$ is indeed equivalent to the conjunction.

\margin{for another proof of the equivalence of $2^\circ)+3^\circ)$ with $\widetilde\sigma^2=\widetilde\rho^3=1$, cf. lemma below}

Let us note that $2^\circ)$ is also written $\widetilde\sigma^2=1$, and is equivalent to: $\unsure{w|\omega}'=0$, i.e. $\unsure{w|\omega}=1$, and we then find, writing \deleted{\unsure{$\rho=$}} 
\[
\rho=\begin{pmatrix}0&1\\\unsure{v|\nu}&1\end{pmatrix},
\]
\[
(23)\qquad
\rho^3=\begin{pmatrix}\unsure{v|\nu}&\unsure{v|\nu}+1\\\unsure{v|\nu}(\unsure{v|\nu}+1)&2\unsure{v|\nu}+1\end{pmatrix}
\]
thus $\rho^3$ can be a scalar i.e. $\widetilde\rho^3=1$ only if the two diagonal terms are equal, now their difference is $(2\unsure{v|\nu}+1)-\unsure{v|\nu}=\unsure{v|\nu}+1$,

\mspage{326}

\scanpage{328}

so the condition reads $\unsure{v|\nu}=-1$, i.e. this is condition $3^\circ$, but then the expression (23) for $\rho^3$ shows us that $\rho^3=-1$.

\uline{\unsure{Remarks|Remark}} The preceding conditions imply the formulas
\[
\begin{aligned}
(24)\qquad &\varepsilon_0=\sigma\rho\\
(25)\qquad &\varepsilon_1=\rho\sigma\\
(26)\qquad &\rho\varepsilon_1=\varepsilon_0^{-1}\\
(27)\qquad &\rho^2\sigma=\varepsilon_0^{-1}
\end{aligned}
\]
each of which \inserted{except the last (27)} implies them in turn
--- for example we have $\varepsilon_0=-\sigma^{-1}\rho$ \inserted{(7)}, so (24) is equivalent to \inserted{$\sigma=-\sigma^{-1}$ i.e.} $\sigma^2=-1$, and (25) is equivalent to (24) by applying $\operatorname{int}(\rho)$ to both \unsure{sides}.

We have moreover
\[
\rho\varepsilon_1=
\begin{pmatrix}0&1\\\unsure{v|\nu}&\unsure{w|\omega}\end{pmatrix}
\begin{pmatrix}1&0\\1&1\end{pmatrix}
=\begin{pmatrix}1&1\\\unsure{v|\nu}+\unsure{w|\omega}&\unsure{\unsure{w|\omega}^{\prime}|\unsure{w|\omega}}\end{pmatrix},
\qquad
\varepsilon_0^{-1}=\begin{pmatrix}1&1\\0&1\end{pmatrix}
\]
which shows that $\rho\varepsilon_1=\varepsilon_0^{-1}$ is equivalent to:
\[
\unsure{w|\omega}=1,\quad \unsure{v|\nu}+\unsure{w|\omega}=0\qquad\text{i.e.}\quad \unsure{w|\omega}=1,\unsure{v|\nu}=-1\quad\text{i.e.}\quad
\rho=\begin{pmatrix}0&1\\-1&1\end{pmatrix},
\]
OK. Composing (25) \inserted{(in the form $\rho\varepsilon_1=\rho^2\sigma$)} and (26) we find (27) as a consequence of the conditions $2^\circ)+3^\circ)$; \deleted{\unsure{the}}\inserted{Let us try to prove} converse implication \deleted{\unsure{follows}}\inserted{by} \deleted{from} a brute-force calculation
\[\deleted{\rho^2\sigma=}\rho^2\sigma=
\begin{pmatrix}
-\unsure{v|\nu}\unsure{w|\omega}&\unsure{w|\omega}-(\unsure{v|\nu}+\unsure{w|\omega}^2)\\
-\unsure{v|\nu}(\unsure{v|\nu}+\unsure{w|\omega}^2)&\unsure{-\unsure{w|\omega}(\unsure{w|\omega}^2+2\unsure{v|\nu})+\unsure{v|\nu}-\unsure{w|\omega}^2|-\unsure{w|\omega}(\unsure{w|\omega}^2+2\unsure{v|\nu})+\unsure{v|\nu}+\unsure{w|\omega}^2}
\end{pmatrix},
\qquad
\varepsilon_0^{-1}=\begin{pmatrix}1&\unsure{-1/\unsure{v|\nu}}\\0&1\end{pmatrix}
\]
\deleted{since $\unsure{v|\nu}$ invertible} we can have the equality $\rho^2\sigma=\unsure{-\varepsilon_0^{-1}|\varepsilon_0^{-1}}$ only if $-\unsure{v|\nu}(\unsure{v|\nu}+\unsure{w|\omega}^2)=0$, \unsure{this implies|which implies} ($\unsure{v|\nu}$ invertible!) $\unsure{v|\nu}+\unsure{w|\omega}^2=0$, \deleted{\unsure{whence}}\inserted{i.e. $\unsure{v|\nu}=-\unsure{w|\omega}^2$},

\mspage{327}

\scanpage{329}
\mspage{465}

\[
\rho^2\sigma=\begin{pmatrix}\unsure{w|\omega}^3&\unsure{w|\omega}\\0&\deleted{\unsure{-w(1+2w)}}\inserted{\unsure{w|\omega}^3}\end{pmatrix},
\]
This shows that
\deleted{$\rho^2\sigma=\varepsilon_0^{-1}$ \unsure{iff} $\unsure{w|\omega}=1$ OK.}
\margin{\deleted{\unsure{We have in addition to} \illegible{4}}
\unsure{$\unsure{v|\nu}=-\unsure{w|\omega}^2$} \deleted{\illegible{2}}
\unsure{i.e. $\unsure{v|\nu}=-1,\unsure{w|\omega}=1$}}
\inserted{
\[
(27\mathrm{bis})\qquad
\rho^2\sigma=\varepsilon_0^{-1}\Longleftrightarrow
\left\{\unsure{v|\nu}+\unsure{w|\omega}^2=0,\ \unsure{w|\omega}^3=1\right.
\quad
\left(\Rightarrow \unsure{v|\nu}=-\unsure{w|\omega}^2=-1/\unsure{w|\omega}\ \text{i.e.}\ \unsure{w|\omega}=-1/\unsure{v|\nu},\ \mathrm{OK}\right)
\]
}

But (27), on multiplying by $\rho$, is also equivalent to
\[\deleted{\unsure{(27)}\,\illegible{1}}\rho^3\sigma=\deleted{\unsure{u}}\rho\varepsilon_0^{-1}
\]
now we know that $\varepsilon_0=\unsure{-\sigma^{-1}\rho}$ \inserted{(7)} i.e. $\varepsilon_0^{-1}=-\deleted{\unsure{\sigma}}\rho^{-1}\unsure{\sigma|\sigma_0}$
i.e. $\rho\varepsilon_0^{-1}=-\sigma$, so the preceding relation also reads $\rho^3\sigma=-\sigma$ or again
\[
(28)\qquad \boxed{\rho^3=-1}.
\]
We can see \deleted{\unsure{this}} the equivalence \deleted{directly}\inserted{of (28) with the conditions (27bis)} by computing $\rho^3$ in terms of $\unsure{v|\nu},\unsure{w|\omega}$, we find
\[
(29)\qquad
\rho^3=\begin{pmatrix}\unsure{v|\nu}\unsure{w|\omega}&\unsure{v|\nu}+\unsure{w|\omega}^2\\\unsure{v|\nu}(\unsure{v|\nu}+\unsure{w|\omega}^2)&\unsure{w|\omega}(2\unsure{v|\nu}+\unsure{w|\omega}^2)\end{pmatrix}
\]
And we can have $\widetilde\rho^3=1$ i.e. $\rho^3$ a scalar only if \deleted{$\unsure{v|\nu}+\unsure{w|\omega}^2=$}
\[
\unsure{v|\nu}+\unsure{w|\omega}^2=0\qquad\text{i.e.}\quad \unsure{v|\nu}=-\unsure{w|\omega}^2
\]
which implies
\[
\rho^3=\begin{pmatrix}-\unsure{w|\omega}^3&0\\0&-\unsure{w|\omega}^3\end{pmatrix}
\]
i.e. we find

\uline{Lemma} We have $\widetilde\rho^3=1$ iff $\unsure{v|\nu}+\unsure{w|\omega}^2=0$ i.e. $\unsure{v|\nu}=-\unsure{w|\omega}^2$ i.e. we have
\[
(30)\qquad
\rho=\begin{pmatrix}0&1\\-\unsure{w|\omega}^2&\unsure{w|\omega}\end{pmatrix},
\]
and then we have
\[
(31)\qquad \rho^3=\unsure{-}\unsure{w|\omega}^3\cdot1\ ,
\]
so we have
\[
(32)\qquad
\rho^3=-1\Longleftrightarrow
\deleted{\unsure{\unsure{v|\nu}=-\unsure{w|\omega}^2}}\quad
\begin{gathered}
\unsure{w|\omega}^3=1\ \text{i.e. }\unsure{w|\omega}\in\Gamma\mu_3\\
\inserted{\uline{and}\quad \unsure{v|\nu}=-\unsure{w|\omega}^2\ (= -1/\unsure{w|\omega})}
\end{gathered}
\]

\mspage{328}

\scanpage{330}

\uline{Corollary 1} The equivalent conditions (22) are also equivalent to
\[
(33)\qquad \rho^3=-1\quad\text{and}\quad\det\rho=1
\]
or again to
\[
(34)\qquad \rho^3=-1\quad\text{and}\quad\deleted{\unsure{\det}}\inserted{\operatorname{Tr}}\rho=1
\]
in other words, considering the three relations
\[
(35)\qquad \det\rho=1,\quad\operatorname{Tr}\rho=1,\quad\rho^3=-1
\]
\inserted{the conjunction of} two of these \unsure{imply|implies} the third, and is equivalent to the equivalent conditions of (22).

\margin{false}
\deleted{
Let us note another way of finding this corollary, by noting that
\[
(\unsure{36|30})\qquad
\rho^3+1=(\rho+1)(\rho^2-\rho+1)\qquad
\text{(algebraic \uline{identity})}
\]
now $\rho$ not being a scalar on any fiber,
\uline{$\rho+1$ is invertible}, so $\rho^3=-1$ i.e. $\rho^3+1=0$ is equivalent to the relation $\rho^2-\rho+1=0$:
\[
(36)\qquad \rho^2-\rho+1=0\ ;\quad(\Longleftrightarrow\rho^3=-1),
\]
which we compare with Hamilton--Cayley
\[
(37)\qquad
\rho^2-\underbrace{(\operatorname{Tr}\rho)}_{\unsure{w|\omega}}\rho+\underbrace{\det\rho}_{-\unsure{v|\nu}}=0,
\]
so \deleted{$\unsure{w|\omega}=1,\unsure{v|\nu}=-1$} $\operatorname{Tr}\rho=\det\rho=1$ implies (36), \deleted{so $\rho^3=-1$}
on the other hand if we have (36), we conclude via 37
\[
(\operatorname{Tr}(\rho)-1)\rho\ \unsure{=|\ne}\ \deleted{\unsure{\det}}\,(\det\rho-1)1
\]
}

\mspage{329}

\scanpage{331}
\mspage{466}

Another way of finding certain equivalences is via the Hamilton--Cayley equations for $\sigma$ and $\rho$
\[
(36)\qquad
\deleted{\unsure{\rho}}\inserted{\sigma}^2
-\underbrace{(\operatorname{Tr}\sigma)}_{\unsure{w|\omega}'=1-\unsure{w|\omega}}\sigma
+\underbrace{\det\sigma}_{\unsure{v|\nu}'=-\unsure{v|\nu}}=0
\]
\[
(37)\qquad
\rho^2-\underbrace{(\operatorname{Tr}\rho)}_\unsure{w|\omega}\rho
+\underbrace{\det\rho}_{-\unsure{v|\nu}}=0.
\]
Equation (36) gives again that the conjunction of $\operatorname{Tr}\sigma=0$, $\det\sigma=1$ implies $\sigma^2+1=0$ i.e. $\sigma^2=-1$, and conversely, taking into account that $1$ and $\sigma$ are \inserted{sections} \unsure{linear\textsuperscript{ly}} independent \unsure{in|of} $\unsure{\mathrm{M}_2}(\unsure{\underline{O}_S|\mathcal O_S})$.
\unsure{Likewise|=} (37) shows
\[
(38)\qquad
(\operatorname{Tr}\rho=1\ \text{and}\ \det\rho=1)
\Longleftrightarrow(\rho^2-\rho+1=0).
\]
and we \deleted{relation} have the implication
\[
(39)\qquad
\rho^2-\rho+1=0\quad\Longrightarrow\quad
\rho^3+1=0\ \text{i.e.}\ \rho^3=-1
\]
by virtue of the \uline{algebraic} identity
\[
\rho^3+1=(\rho+1)(\rho^2-\rho+1).
\]
So

\uline{Corollary 2}. The equivalent conditions of (22) are also equivalent to the relation
\[
(40)\qquad \rho^2+\rho+1=0.
\]

\mspage{330}

\scanpage{332}

\deleted{
$5^\circ)$ Modular pinning theorem.

\uline{Theorem} Let $G,SG,PG$ be forms \inserted{(over scheme $S$)} of $\mathrm{Gl}(2),\mathrm{Sl}(2),\mathrm{GP}(1)$ \inserted{\deleted{\unsure{over scheme}}} \unsure{(}associated\unsure{)}, and $\widetilde\sigma,\widetilde\rho$ two sections of $PG$.
}

$5^\circ)$ \uline{A return to} $\uline{1^\circ)}$, concerning the situation
\[
(41)\quad
\left\{\begin{array}{l}
\widetilde\rho,\widetilde\sigma\ \text{with }1^\circ)\quad
\widetilde\varepsilon_0=\widetilde\sigma^{-1}\widetilde\rho\quad
\uline{\text{regular}}\text{ unipotent}\ \deleted{\text{\unsure{i.e.}}}\quad(\text{i.e.}\\
\widetilde\varepsilon_1=\widetilde\rho\widetilde\sigma^{-1}
=\widetilde\rho(\widetilde\varepsilon_0)
=\widetilde\sigma(\widetilde\varepsilon_0)\quad
\uline{\text{regular}}\text{ unipotent})\ \text{and}\\
2^\circ)\quad\widetilde\varepsilon_0\ \text{and}\ \widetilde\varepsilon_1\ \text{inequivalent in each fiber.}
\end{array}\right.
\]

Let $\varepsilon_0,\varepsilon_1$ be the \uline{regular} unipotent liftings in $G$ --- identified with sections of $\unsure{M|\mathcal M}$
\[
(42)\qquad\varepsilon_0=1+N_0,\quad\varepsilon_1=1+N_1,
\]
\margin{$N_0$ and $N_1$ strictly \unsure{nilpoten\textsuperscript{t}} i.e. their determinant and their trace \deleted{\illegible{1}} zero}
with
\[
(43)\qquad
\widetilde\sigma(\varepsilon_0)=\widetilde\rho(\varepsilon_0)=\varepsilon_1,
\quad\inserted{\text{i.e.,}}\quad
\widetilde\sigma(N_0)=\widetilde\rho(N_0)=N_1.
\]
We can then lift $\widetilde\sigma$ and $\widetilde\rho$ canonically to sections $\sigma,\rho$ of $G$ (or of $\unsure{M|\mathcal M}$) satisfying the following conditions
\[
(44)\quad\left\{
\begin{array}{ll}
1^\circ)\ -\sigma^{-1}\rho=\varepsilon_0
&[\Longleftrightarrow-\rho\sigma^{-1}=\varepsilon_1]\\
\deleted{\unsure{2^\circ)(\rho+\sigma)}}
&\inserted{\left[\begin{aligned}
&\Longleftrightarrow\rho+\sigma=-N_1\sigma\ (= -\sigma N_0)\\
&\Longleftrightarrow\rho+\sigma=N_1\rho\ (=\rho N_0)
\end{aligned}\right]}\\
&\deleted{\unsure{\Longleftrightarrow\rho+\sigma=}}\\
2^\circ)\ (\rho+\sigma)^2=\rho+\sigma
&\inserted{\begin{gathered}
\Longrightarrow(\rho+\sigma)^2=\unsure{u}(\rho+\sigma)\quad\text{with}\\
u\in\Gamma G_m
\end{gathered}}
\end{array}\right.
\]
Moreover, there exists a \deleted{\unsure{isomorphism}} unique triple $(\varphi,\rho_0,\sigma_0)$
\[
(45)\quad\left\{\begin{array}{l}
G_0=\mathrm{Gl}(2)_S\xrightarrow{\varphi}G\quad(\text{pinning})\\
\rho_0,\sigma_0\in\Gamma(G_0),\quad
\rho_0=\begin{pmatrix}0&1\\\unsure{v|\nu}&\unsure{w|\omega}\end{pmatrix},\quad
\sigma_0=\begin{pmatrix}0&-1\\\unsure{v|\nu}'&\unsure{w|\omega}'\end{pmatrix}\\
(\unsure{v|\nu}+\unsure{v|\nu}'=0,\unsure{w|\omega}+\unsure{w|\omega}'=1),\quad \unsure{v|\nu}\in\Gamma G_m\\
\inserted{\text{i.e. }\unsure{v|\nu}'=-\unsure{v|\nu},\unsure{w|\omega}'=1-\unsure{w|\omega}}
\end{array}\right.
\]
such that we have
\[
(\unsure{46|48})\qquad\varphi(\rho_0)=\rho,\quad\varphi(\sigma_0)=\sigma,
\]
which implies
\[
(47)\qquad\varphi((\varepsilon_0)_0)=\varepsilon_0,\quad
\varphi((\varepsilon_1)_0)=\varepsilon_1
\]

\mspage{331}

\scanpage{333}
\mspage{467}

where \deleted{\illegible{1}}
\[
(48)\qquad
\unsure{(\varepsilon_0)_0}=-\sigma_0^{-1}\rho_0=
\begin{pmatrix}1&\dfrac{1}{\unsure{v|\nu}}\\0&1\end{pmatrix},\qquad
\deleted{\unsure{(\varepsilon_1)_0}}
\inserted{(\varepsilon_1)_0=-\rho_0\sigma_0^{-1}}=
\begin{pmatrix}1&0\\1&1\end{pmatrix}
\]
--- and conversely, for $\rho_0,\sigma_0$ as above, and any isomorphism
$\varphi:G_0=\mathrm{Gl}(2)_S\xrightarrow{\sim}G$, the images
$\widetilde\rho,\widetilde\sigma$ of $\rho,\sigma$ in $PG$ satisfy
$(41)$, and the images $\rho,\sigma$ of $\rho_0,\sigma_0$ under $\varphi$
are liftings of $\widetilde\rho,\widetilde\sigma$ satisfying $(44)$.

\uline{Modular pinning theorem.} Let $G,SG,PG$ be forms
(\uuline{over} scheme $S$) of $\mathrm{Gl}(2),\mathrm{Sl}(2),\mathrm{GP}(1)$,
\unsure{associated} with an Azumaya \unsure{algebra} $M$ of rank $4$.

$1.)$ Let us call a \uline{modular \deleted{\illegible{1}}\inserted{triple}}
in $PG$ a triple of sections $(\widetilde\rho,\widetilde\sigma,
\widetilde\varepsilon_0)$ of $PG$, satisfying conditions $(49)(50)(51)$.

\deleted{$(\unsure{49})\quad \widetilde\rho^3=\widetilde\sigma^2=1$, $\widetilde\varepsilon_0$ is \unsure{regular} unipotent.}

\[
(49)\qquad \widetilde\varepsilon_0=\widetilde\sigma^{-1}\widetilde\rho
\quad\text{i.e.}\quad
\widetilde\rho^{-1}\widetilde\sigma\widetilde\varepsilon_0=1
\]
[so that
\[
(49\ \mathrm{bis})\qquad
\widetilde\varepsilon_1\overset{\mathrm{dfn}}=
\widetilde\rho\widetilde\sigma^{-1}=
\widetilde\rho(\widetilde\varepsilon_0)=
\widetilde\sigma(\widetilde\varepsilon_0)
\]
];

$(50)\qquad \widetilde\rho^3=\widetilde\sigma^2=1$,\quad
$\widetilde\varepsilon_0$ is regular unipotent
($\Longleftrightarrow\widetilde\varepsilon_1$ is \uline{regular} unipotent);

$(51)\qquad \widetilde\varepsilon_0$ and $\widetilde\varepsilon_1$
\deleted{\unsure{everywhere}} inequivalent \deleted{\unsure{on the}}
\inserted{in each} fiber
(which means, introducing the liftings
\[
\varepsilon_0=1+N_0,\qquad \varepsilon_1=1+N_1
\]
\margin{($\unsure{N_0\text{ and }N_1}$ strictly nilpotent,\ldots\ cf. $(42)$)}
to $G$, that $N_0$ and $N_1$ not proportional in every fiber).

$2)$ Let us call a modular triple \unsure{of|in} $G$ a
\mspage{332}

\scanpage{334}

triple of sections $(\rho,\sigma,\varepsilon_0)$ satisfying the
\deleted{\unsure{five}} following conditions

\[
\begin{aligned}
(52)\qquad &-\sigma^{-1}\rho=\varepsilon_0
\quad(\Longleftrightarrow\ \rho^{-1}\sigma\varepsilon_0=-1.)\\
&\left[\begin{aligned}
&\Longleftrightarrow\ -\rho\sigma^{-1}=\varepsilon_1
&&\text{where }\varepsilon_1=\rho\varepsilon_0\rho^{-1}\\
&\Longleftrightarrow\ -\rho\sigma^{-1}=\varepsilon_1
&&\text{where }\varepsilon_1=\sigma\varepsilon_0\sigma^{-1}
\end{aligned}\right]
\end{aligned}
\]
\[
(53)\qquad(\rho+\sigma)^2=\rho+\sigma
\]

$(54)\qquad\varepsilon_0$
\deleted{and $\varepsilon_1\overset{\mathrm{dfn}}=\rho\varepsilon_0\rho^{-1}=\sigma\varepsilon_0\sigma^{-1}$ \unsure{are}}
\inserted{is regular unipotent i.e. $\varepsilon_0=1+N_0$ with}
$N_0$ \deleted{\unsure{unipotent}} regular nilpotent
($\mathrm{Tr}\,N_0=\det N_0=0$, $N_0$ nonzero in each fiber).

\margin{$\Longleftrightarrow\left\{\begin{array}{l}
\det\varepsilon_0=1\\
\mathrm{Tr}\,\varepsilon_0=2\\
\varepsilon_0\ne1\text{ in every fiber}
\end{array}\right.$}

($\Longleftrightarrow\ \deleted{\unsure{\varepsilon_0}}\ \varepsilon_1
\overset{\mathrm{dfn}}=\rho\varepsilon_0\rho^{-1}
=\sigma\varepsilon_0\sigma^{-1}$ is regular unipotent)

$(55)\qquad\deleted{\rho}\cdot\varepsilon_0$ and $\varepsilon_1$ are
inequivalent in each fiber
($\Longleftrightarrow\ N_0$ and $N_1$ not proportional in every fiber).

$(56)\qquad$ We have one of the four equivalent conditions
\[
\begin{array}{cc}
\boxed{\sigma^2+1=0}&\boxed{\rho^2\unsure{+|-}\rho+1=0}\\[6pt]
\boxed{\mathrm{Tr}\,\sigma=0,\det\sigma=1}&
\boxed{\mathrm{Tr}\,\rho=\det\rho=1}
\end{array}
\quad.
\]

\deleted{\unsure{3}} $3)$ Let us call the \uline{\unsure{standard}}
\uline{modular triple} in $\mathrm{Gl}(2)_S$
(resp. in $\mathrm{GP}(1)_S$) the triple
\[
(57)\quad\left(\rho_0=\begin{pmatrix}0&1\\-1&1\end{pmatrix},\quad
\sigma_0=\begin{pmatrix}0&-1\\1&0\end{pmatrix},\quad
\varepsilon_0=\begin{pmatrix}1&-1\\0&1\end{pmatrix}\right)
\]
(resp. its image $(\widetilde\rho_0,\widetilde\sigma_0,
\unsure{\widetilde\varepsilon_0})$ in $\Gamma\mathrm{GP}(1)_S$) ---
which is indeed a modular triple.

This being so, we have \deleted{some}\inserted{the \unsure{commutative} diagram of} canonical
maps
\mspage{333}

\scanpage{335}
\mspage{468}

\[
\varphi\longmapsto(\varphi(\rho_0),\varphi(\sigma_0),\varphi(\varepsilon_0))
\]

$(58)$
\diagram{\begin{tikzcd}[ampersand replacement=\&,column sep=large,row sep=large]
\operatorname{Isom}(\mathrm{Gl}(2),G)
\arrow[r]\arrow[d,phantom,"\unsure{\wr}"]\&
\mathrm{TripMod}(G)\arrow[d]\&
(\rho,\sigma,\varepsilon_0)\arrow[d,mapsto]\\
\operatorname{Isom}(\mathrm{GP}(1),PG)\arrow[r]\&
\mathrm{Tripmod}(PG)\&
(\widetilde\rho,\widetilde\sigma,\widetilde\varepsilon_0)
\end{tikzcd}}
\[
\psi\longmapsto(\psi(\widetilde\rho_0),\psi(\widetilde\sigma_0),
\psi(\widetilde\varepsilon_0))
\]

and we have the

\uline{Pinning theorem.} All the arrows of $(58)$ are bijections.

This theorem only summarizes \deleted{the}\inserted{some}
\deleted{\unsure{reflections}}\inserted{of the} main results of the reflections
\unsure{of the|of the} present \S\ \uline{$\overline{\mathrm{VI}}$}.
Let us note that the injectivity of the horizontal arrows means \inserted{also}
that an \deleted{\unsure{isomorphism}}
$\mathrm{Gl}(2)_S=G_0\longrightarrow G$
(resp. $\mathrm{GP}(1)_S=PG_0\longrightarrow\unsure{G|PG}$
is known when we know its effect on $\sigma$ and $\rho$
(thus also on $\varepsilon_0$]
\unsure{$\leftarrow$} resp. when we know its effect on
$\widetilde\sigma$ and $\widetilde\rho$) --- which also means
the \unsure{same|=} thing for the automorphisms of
$\mathrm{Gl}(2)_S$ (resp. of $\mathrm{GP}(1)_S$)
which also means that
\[
(59)\qquad T_\rho\cap T_\sigma=\deleted{\unsure{\mu_2|\{1\}}}\ G_m
\]
(the intersection of the centralizers of $\rho,\sigma$ in
$\mathrm{Gl}(2)_S$ is \unsure{reduced|reduced} to the center $G_m$) ---
\unsure{resp.} the stronger relation
\[
(60)\qquad N'_\rho\cap N'_\sigma=G_m,
\]
\mspage{334}

\scanpage{336}

where $N'_\rho$, $N'_\sigma$ are the group \unsure{subschemes} of $\mathrm{Gl}(2)_S$ defined \deleted{by} \inserted{\unsure{respective\textsuperscript{ly}}} as stabilizers of $\widetilde\rho$, $\widetilde\sigma$ in $\unsure{\mathrm{Gl}(2)_S}$, i.e.
\[
(61)\qquad
\left\{\begin{aligned}
N'_\rho&=\{g\mid g(\rho)=\varepsilon\rho,\ \varepsilon\text{ scalar}\},\\
N'_\sigma&=\{g\mid g(\sigma)=\varepsilon\sigma,\ \varepsilon\text{ scalar}\}
\end{aligned}\right.
\]
\margin{N.B. we \unsure{will have} \unsure{necessari\textsuperscript{ly}} $\varepsilon^2=1$}

\uline{Remark} It is plausible, \deleted{that \illegible{5}} \inserted{\unsure{In view of these results,}} that $\mathrm{GP}(1)_{\mathbb Z}$ --- as an ``absolute'' sheaf of groups (over $\operatorname{Spec}\mathbb Z$) for fppf, is generated by the homomorphisms
\[
\begin{aligned}
\mathbb Z/2\mathbb Z&\xrightarrow{\alpha\,\mapsto\,\widetilde\sigma_0^{\alpha}}\mathrm{GP}(1)_{\mathbb Z}\\
\mathbb Z/3\mathbb Z&\xrightarrow{\beta\,\mapsto\,\widetilde\rho_0^{\beta}}\mathrm{GP}(1)_{\mathbb Z}\\
\mathbb G_a&\xrightarrow{\lambda\,\mapsto\,\widetilde\varepsilon_0(\lambda)}\mathrm{GP}(1)_{\mathbb Z}
\end{aligned}
\]
satisfying \unsure{the|the} \inserted{sole} relation (\unsure{R})
\[
(62)\qquad \widetilde\rho_0^{-1}\widetilde\sigma_0\widetilde\varepsilon_0(1)=1
\qquad\deleted{\unsure{(?)}}
\]
--- it would be well worth the trouble to clear up this question, as well as the analogous question for $\mathrm{Gl}(2)_{\mathbb Z}$, where the \deleted{relations} groups $\mathbb Z/2\mathbb Z$, $\mathbb Z/3\mathbb Z$ must \unsure{be replaced} by $\mathbb Z/4\mathbb Z$, $\mathbb Z/6\mathbb Z$, and the relation (62), by
\[
(63)\qquad \rho_0^{-1}\sigma_0\widetilde\varepsilon_0(1)=\rho_0^3=\sigma_0^2.
\]
\mspage{335}

\scanpage{337}
\mspage{469}

\uline{Remark} I have the distinct impression that the package of conditions (51) to (56) is redundant, and wonder for example whether \unsure{the set|one} of the three simple algebraic conditions
\[
(64)\qquad \sigma^2+1=\rho^2-\rho+1=0,\qquad (\rho+\sigma)^2=\rho+\sigma
\]
does not \unsure{imply} the other conditions \inserted{(defining $\varepsilon_0$ by (\unsure{51|52}))}. The third condition (64) means that
\[
(65)\qquad N=\rho+\sigma
\]
is \unsure{idempotent}
\[
(66)\qquad N^2=N,
\]
and in terms of $N$ we will have
\[
\sigma^{-1}N=\sigma^{-1}\rho+1\quad\text{i.e.}
\]
\[
(67)\qquad
\left\{\begin{aligned}
\varepsilon_0&\overset{\mathrm{dfn}}{=}-\sigma^{-1}\rho=1-\sigma^{-1}N=1+N_0
&\qquad N_0&=-\sigma^{-1}N\\[-2pt]
&&&=\sigma N\\
\varepsilon_1&\overset{\mathrm{dfn}}{=}-\rho\sigma^{-1}=1-N\sigma^{-1}=1+N_1
& N_1&=-N\sigma^{-1}=N\sigma
\end{aligned}\right.
\]
\margin{\uline{N.B.} We must \deleted{\unsure{perhaps}} \uline{specify} \deleted{\unsure{that in}} \illegible{4} \unsure{first} relations \unsure{(64)} \illegible{3}
$\operatorname{Tr}\sigma=0$, $\det\sigma=1$, $\operatorname{Tr}\rho=\unsure{0|1}$, $\det\rho=1$ \unsure{must require that} \illegible{2} are \uline{regular} i.e. \unsure{not homotheties} in any fiber \illegible{1}.}

and the condition $\varepsilon_0$ \deleted{\unsure{strict\textsuperscript{ly} unipotent}} \inserted{\unsure{regular unipotent}} is expressed by the conditions
\[
(68)\qquad \det\sigma N=\operatorname{Tr}\sigma N=0,\quad
\begin{gathered}
\sigma N\text{ nonzero in every fiber.}\\
\text{i.e. }N\text{ nonzero in every fiber}
\end{gathered}
\]
Now the \deleted{conditions} \inserted{third of these} (68) follows from the others, for if we had $N=0$ on a fiber, we could suppose $S$ the spectrum of a field, $N=0$, i.e. $\rho=-\sigma$, so \deleted{\illegible{1}} we would have both
\[
\sigma^2+1=0\qquad \sigma^2+\sigma+1=0
\qquad\text{whence }\sigma=0,\text{ absurd.}
\]
I also say that in no fiber do we have $N=1$, otherwise we would have in this fiber \deleted{$\rho=1-\sigma$} $\sigma=1-\rho$ and substituting in $\sigma^2+1=0$ we find \deleted{\illegible{1}}
\[
(1-\rho)^2+1=0\quad\text{i.e.}\quad \rho^2-2\rho+2=0,
\]
\mspage{336}

\scanpage{338}

comparing with $\rho^2-\rho+1=0$, we find \deleted{$\unsure{\rho-1}$} by subtraction $\rho-1=0$ i.e. $\rho=1$, but the relation $\rho^2-\rho+1=0$ becomes $1=0$, absurd.

Thus the idempotent operator $N$ is \deleted{everywhere} of rank $1$ \inserted{in each fiber}, so \illegible{1} (unless I am mistaken) \deleted{\unsure{says}} that $N$ is \deleted{decomposable} (in terms of $\unsure{\mathcal M}\simeq\check F\otimes F$) \inserted{local\textsuperscript{ly}} decomposable in the form $x'\otimes x$, with $\langle x,x'\rangle=1$, and \deleted{$\operatorname{Tr}$}
\[
\operatorname{Tr}N=1,\quad\det N=0
\]
so the relation $\det(\sigma N)=0$ in (68) is also automatic, since $\det\sigma N=\det\sigma\det N=0$. There remains only the relation
\[
\operatorname{Tr}\underbrace{\sigma N}_{\substack{\sigma^2+\sigma\rho\\=\sigma\rho-1}}=0
\]

( \deleted{that} which is also written
\[
(69)\qquad \operatorname{Tr}\sigma\rho=2\quad(\Longleftrightarrow\operatorname{Tr}\sigma^{-1}\rho=-2)
\]
\unsure{for} which we would have to see whether or not it follows from the other relations.

Let us take
\[
(70)\qquad
\left\{\begin{aligned}
\rho&=\rho_0=\begin{pmatrix}0&1\\-1&1\end{pmatrix}\\
\sigma&=\underline u\sigma_0\underline u^{-1}
=\begin{pmatrix}\unsure{a}&b\\c&d\end{pmatrix}
\begin{pmatrix}0&-1\\1&0\end{pmatrix}
\begin{pmatrix}d&-b\\-c&a\end{pmatrix}\\
&\text{where }\underline u=\begin{pmatrix}a&b\\c&d\end{pmatrix}\in\unsure{\Gamma\mathrm{Sl}(2)},\text{ i.e. }ad-bc=1
\end{aligned}\right.
\]
We find
\mspage{337}

\scanpage{339}
\mspage{470}

\[
(70\mathrm{bis})\qquad
\sigma=\begin{pmatrix}
bd+ac&-(a^2+b^2)\\
c^2+d^2&-(bd+ac)
\end{pmatrix}
\]
whence
\[
N=\rho+\sigma=\begin{pmatrix}
(bd+ac)&1-(a^2+b^2)\\
-1+(c^2+d^2)&1-(bd+ac)
\end{pmatrix}
\]
\[
N-1=\begin{pmatrix}
(bd+ac)-1&1-(a^2+b^2)\\
-1+(c^2+d^2)&-(bd+ac)
\end{pmatrix}
\quad\begin{gathered}
=-\det(N)N^{-1}\\
\text{if }\det N\text{ \unsure{invertible}}
\end{gathered}
\]
and we find
\[
(71)\qquad N^2-N=N(N-1)=-\det(N).1
\]
so \deleted{the condition}
\[
(72)\qquad N^2=N\deleted{\unsure{\Longleftrightarrow}}\quad\Longleftrightarrow\quad\det N=0\ ;\quad(\Longrightarrow\operatorname{Tr}N=1)
\]
Now we have
\[
\begin{aligned}
\det N={}&bd+ac-(bd\unsure{+}ac)^2\deleted{\unsure{+1}}\\
&\unsure{-|=}\bigl[-1-(a^2+b^2)(c^2+d^2)+a^2+b^2+c^2+d^2\bigr]\\
={}&bd+ac-\deleted{b^2d^2}-\deleted{a^2c^2}-2a\uline{bc}d+1
+\deleted{a^2c^2}+\uline{a^2d^2}+\uline{b^2c^2}+\deleted{\unsure{b^2d^2}}\\
&\hspace{4em}-a^2-b^2-c^2-d^2\\
={}&(bd+ac)-(a^2+b^2+c^2+d^2)+\underset{1}{\underbrace{(ad-bc)^2}}+1\\
={}&2+(bd+ac)-(a^2+b^2+c^2+d^2)
\end{aligned}
\]
so the condition \unsure{(72)|(71)} is equivalent to
\[
(73)\qquad a^2+b^2+c^2+d^2=2+(bd+ac),
\]
on the other hand we have by 70, (70bis)
\[
\unsure{\rho\sigma}=\begin{pmatrix}
c^2+d^2&-(bd+ac)\\
c^2+d^2-(bd+ac)&a^2+b^2-(bd+ac)
\end{pmatrix}
\]
so
\[
(74)\qquad \operatorname{Tr}\unsure{\sigma\rho}=\operatorname{Tr}\rho\sigma
=a^2+b^2+c^2+d^2-(bd+ac)=2
\]
\mspage{338}
